\section{A semantic--structural factorial benchmark}

The central benchmark independently manipulates surface-semantic similarity and structural similarity. Let $S_{\mathrm{sem}}\in\{\mathrm{low},\mathrm{high}\}$ and $S_{\mathrm{str}}\in\{\mathrm{low},\mathrm{high}\}$. The four cells are:

\begin{center}
\begin{tabular}{lll}
\toprule
Cell & Semantic relation & Structural relation \\
\midrule
Q1 & high & high \\
Q2 & low & high \\
Q3 & high & low \\
Q4 & low & low \\
\bottomrule
\end{tabular}
\end{center}

Q2 and Q3 are the load-bearing cells. Q2 asks whether the system can transfer when vocabulary and domain semantics are deliberately distant but the relevant relation/invariant structure is preserved. Q3 is a semantic decoy: vocabulary is similar and embedding similarity should be high, but the causal/dynamical structure or boundary condition that justifies transfer is changed. A useful structural system should accept Q2 and reject Q3.

The deterministic scaffold in \path{src/rakl/structural_benchmark.py} now exercises this requirement across three distinct mechanism families rather than one illustrative analogy. For \emph{queue stability}, a computer network and an emergency department use different entity vocabularies but share a continual-arrival/finite-service/backlog structure for the registered QoI; the high-semantic decoy instead uses queue language in a finite-batch regime, so the steady-flow stability invariant is invalid. For \emph{positive feedback}, viral adoption and a bank-run withdrawal loop share the sign structure required for perturbation amplification, while a semantically similar control-feedback target reverses one loop edge and becomes stabilizing negative feedback. For a \emph{threshold cascade}, epidemic spread and cascading grid outage share a supercritical propagation structure, while a same-domain subcritical cascade decoy removes the load-bearing reproduction relation and violates the threshold/connectivity boundary. The original Q4 case remains a semantically and structurally unrelated queue-to-feedback mapping.

\subsection{Deterministic conformance result}

The exact-subject executable receipt contains six load-bearing Q2/Q3 cases: one Q2 and one Q3 case in each of the three mechanism families. The deliberately semantic-only decision rule licenses high-semantic candidates and rejects low-semantic candidates; it is therefore wrong on all six hard cases by construction. The structural gate instead checks the registered directional witness, QoI, load-bearing relations, invariants and target boundaries. It licenses all three Q2 transfers and rejects all three Q3 decoys.

\begin{center}
\begin{tabularx}{\linewidth}{@{}lYYY@{}}
\toprule
Decision rule & Q2 valid licensed & Q3 invalid accepted & Hard cases correct \\
\midrule
Semantic-only control & 0/3 & 3/3 & 0/6 \\
Structural witness gate & 3/3 & 0/3 & 6/6 \\
\bottomrule
\end{tabularx}
\end{center}

This table is a software-conformance result, not a statistical generalization result. The cases were deliberately constructed to separate semantic and structural similarity; their role is to establish that the implementation can express the intended distinction before model-level experiments are authorized. They provide no estimate of performance on naturally occurring cross-domain mappings.

These cases are deterministic conformance worlds, not an empirical claim that the representation generalizes. Their purpose is to prevent a one-family implementation from earning an expensive training run merely by memorizing queue-specific details. The confirmatory dataset will still require independent annotations and broader families such as flow conservation, bottlenecks, delayed negative feedback, diffusion/contagion, matching, scheduling, concentration fragility, path dependence and search/backtracking. Each family should span multiple surface domains and include adversarial near-misses.

\subsection{A broader internal diagnostic fails closed at the annotation gate}

Before inspecting predictor outcomes, we froze a balanced Q1--Q4 proposal set, leave-one-family-out split, predictor family, metrics and stopping thresholds at Git subject \texttt{f2701f732f83}. The resulting diagnostic contains 44 constructed proposal pairs across 11 proposed mechanism families. Every coordinate was proposed in the same assistant session that assembled the benchmark. None has two independent human/expert annotations or adjudication, so all 44 are explicitly marked ineligible for confirmatory use.

The local pilot fits deterministic L2-regularized logistic predictors using successively richer feature sets: calibrated surface overlap, skill-tag overlap, dependency-pattern overlap and witnessed structure. It is a cheap predictive diagnostic with no foundation-model judgement. Leave-one-family-out outcomes crossed the frozen incremental-signal thresholds: relative to the strongest non-witness control, the dependency-aware model, witnessed structure increased ROC-AUC from 0.914 to 1.000 and average precision from 0.903 to 1.000, while reducing Brier loss from 0.166 to 0.039. Its Q2 true-accept rate was 1.00 and Q3 false-accept rate was 0.00. These point estimates measure discrimination of labels encoded in an internally constructed proposal set; they do not estimate human agreement or natural transfer performance.

\begin{figure}[t]
\centering
\includegraphics[width=\textwidth]{figures/paper3_cheap_gate_internal.pdf}
\caption{Receipt-bound internal cheap-gate diagnostic. The unit is a constructed source--target pair ($n=44$) and each prediction is out-of-family under a deterministic leave-one-family-out fit. Bars are aggregate point estimates; there are no repeated seeds, sampling intervals or independent annotations. Panel a reports ROC-AUC and average precision, panel b Brier and log loss, and panel c Q2 true accept with Q3 false accept. The figure is generated only from the checked machine-readable result receipt; values are not hand-entered into the plot.}
\label{fig:cheap-gate-internal}
\end{figure}

The independent-annotation gate nevertheless has authority over the diagnostic-signal gate. Its exact verdict is \texttt{FAIL\_CLOSED\_MISSING\_INDEPENDENT\_ANNOTATION}: zero of 44 items are confirmatory-eligible. Consequently, no training or inference run was launched. The unresolved residual is evidence acquisition, not a license to substitute same-session agreement or spend compute defending the original shared-substrate headline.

A subsequent chronology audit identified a stricter cut: the frozen v1 protocol records both that labels were visible during construction and that confirmatory use is not permitted.  The 44-item v1 set is therefore permanently an internal constructed diagnostic; later annotations cannot retroactively convert it into confirmatory evidence or authorize compute.  We retain that receipt unchanged as negative history and register a separate v2 lane.  V2 requires a fresh label-blind item set frozen before labels or results, the same or stricter signal thresholds, two genuinely independent human/expert annotations per item, distinct post-freeze adjudication and an external provenance audit.  Repository validation can check hashes, coverage, chronology and attestations, but cannot itself prove that an independence attestation is truthful.  Until both the new annotation gate and held-out signal gate pass, the LUNARC preflight refuses training and inference submission.

\subsection{Structural score must add information beyond semantic score}

For a transfer outcome $Y$, the confirmatory analysis compares models such as
\[
\operatorname{logit}P(Y=1)=\beta_0+\beta_1 S_{\mathrm{sem}}
\]
against
\[
\operatorname{logit}P(Y=1)=\beta_0+\beta_1 S_{\mathrm{sem}}+\beta_2 S_{\mathrm{str}}+\beta_3S_{\mathrm{sem}}S_{\mathrm{str}}+u_{\mathrm{family}}.
\]
The central mechanism requires material incremental predictive value from the structural term and, operationally, a low invalid-transfer rate in Q3. If a strong semantic retriever or LLM analogy judge already predicts valid transfer as well as the witnessed structure, the proposed representation does not earn an additional layer.

\subsection{Annotation and leakage}

A human-curated version will record role graphs, typed relations, invariants, boundary conditions and QoI with at least two independent annotations plus adjudication. Inter-annotator agreement is reported separately for roles, relations, class assignment and boundaries. The final domain split must prevent a surface domain from appearing in both training and target evaluation for the same structure family. Synthetic generation can expand coverage, but held-out human adjudication is required to show that the structural labels are not merely artifacts of the generator.
